% ============================================================
% 00-sazetak-hr.tex
% ============================================================

Problem neuravnote\v{z}enosti klasa jedan je od klju\v{c}nih izazova u
strojinom u\v{c}enju, posebice u klasifikacijskim zadacima gdje
distribucija primjera me\dj{}u klasama nije uniformna. Algoritam SMOTE
(engl.\ Synthetic Minority Over-sampling Technique) predstavlja jedan od
najpopularnijih pristupa za rje\v{s}avanje ovog problema — generira
sinteti\v{c}ke primjere manjinske klase linearnom interpolacijom
postoje\'{c}ih.

U ovom radu analizirani su temeljni SMOTE algoritam i njegove izvedenice:
Borderline-SMOTE, ADASYN, Safe-Level-SMOTE, K-Means SMOTE, SVM-SMOTE,
kombinirani pristupi SMOTE-ENN i SMOTE-Tomek te geometrijska
pro\v{s}irenja G-SMOTE, Random SMOTE i Polynom-Fit SMOTE. Usporedno su
razmatrani i alternativni pristupi poput cost-sensitive u\v{c}enja,
ansambl metoda i primjene neuronskih mre\v{z}a.

Razvijeno programsko rje\v{s}enje implementira navedene algoritme u
Pythonu i omogu\'{c}uje njihovu eksperimentalnu evaluaciju na \v{s}irokom
spektru skupova podataka razli\v{c}itih karakteristika. Provedena je
statisti\v{c}ka analiza rezultata s ciljem utvr\dj{}ivanja pona\v{s}anja
pojedinih algoritama pri razli\v{c}itim omjerima klasa,
dimenzionalnostima i razinama \v{s}uma. Dodatno, izra\dj{}eno je web
su\v{c}elje za vizualizaciju procesa generiranja sinteti\v{c}kih
primjera.

\bigskip
\noindent\textbf{Klju\v{c}ne rije\v{c}i:} klasifikacija, neuravnote\v{z}enost klasa,
preuzorkovanje, sinteti\v{c}ki uzorci, SMOTE
